% Section 1 - The C-to-C++ transition
% Roberto Masocco <roberto.masocco@uniroma2.it>
% April 12, 2026

% ### The C-to-C++ transition ###
\section{The C-to-C++ transition}
\graphicspath{{figs/section1/}}

% --- C and C++ ---
\begin{frame}{C and C++}{What is the relationship?}
	\justifying
	C++ was designed as a \textbg{superset of C}: almost all valid C code is also valid C++ code.\\
	\bigskip
	Everything you know from C is still there:
	\begin{itemize}
		\item primitive types, arrays, and \textbg{pointers};
		\item \textbg{structs}, unions, and enumerations;
		\item operators;
		\item the \textbg{standard C library} (\texttt{<cstdio>}, \texttt{<cstdlib>}, \texttt{<cstring>}...);
		\item manual memory management with \texttt{malloc}/\texttt{free};
		\item \textbg{preprocessor} directives (\texttt{\#include}, \texttt{\#define}, \texttt{\#ifdef}, \texttt{\#endif}...).
	\end{itemize}
	\begin{block}{}
		\centering
		Dust off your C skills. C++ adds \textbf{a lot} on top of them.
	\end{block}
\end{frame}

% --- The C-to-C++ transition ---
\begin{frame}{The C-to-C++ transition}{What changes}
	\justifying
	Some conventions change when moving to C++.\\
	\bigskip
	\textbg{File extensions} reflect the language:
	\begin{itemize}
		\item \texttt{.cpp} (or \texttt{.cc}) for source files;
		\item \texttt{.hpp} (or \texttt{.hh}) for header files.
	\end{itemize}
	\bigskip
	\textbg{Include guards} are replaced by the cleaner pragma:
	\begin{itemize}
		\item \texttt{\#pragma once} at the top of every header file.
	\end{itemize}
	\bigskip
	\textbg{Compilation} uses a C++-aware compiler and specifying the \textbg{language standard} is relevant:
	\begin{itemize}
		\item \texttt{g++ -std=c++17 -o my\_program my\_program.cpp}
	\end{itemize}
	\begin{block}{}
		\centering
		This course targets \textbf{C++17}, the standard required by ROS 2 Jazzy.
	\end{block}
\end{frame}
\begin{frame}{The C-to-C++ transition}{What is added}
	\justifying
	C++ additions are designed to support \textbg{large-scale}, \textbg{maintainable} software; we need to cover:
	\begin{itemize}
		\item \textbg{classes}: encapsulation, constructors, destructors, member functions;
		\item \textbg{RAII}: the most important C++ programming idiom;
		\item \textbg{inheritance} and \textbg{virtual dispatch}: polymorphism with a concrete implementation;
		\item \textbg{references}: a safer, cleaner alternative to pointers in many contexts;
		\item \textbg{namespaces}: organizing and partitioning large codebases;
		\item \textbg{templates}: generic programming, the foundation of the standard library;
		\item \textbg{smart pointers} and the \textbg{standard library}: safer memory and modern containers;
		\item \textbg{lambdas} and \textbg{callbacks}: the backbone of event-driven programming;
		\item \textbg{concurrency}: \texttt{std::thread}, mutexes, condition variables, and atomics.
	\end{itemize}
\end{frame}
\begin{frame}{The C-to-C++ transition}
  \begin{block}{}
    \centering
    This lecture is \href{https://github.com/robmasocco/DAFN26_Robotics_2}{\color{blue}\underline{here}}.\\
    Code examples are available \href{https://github.com/IntelligentSystemsLabUTV/ros2-examples/tree/jazzy/tools/cpp-examples}{\color{blue}\underline{in the \texttt{tools/cpp-examples/} subfolders}} of the materials.\\
    Every lecture ends with \textbf{examples} showing new language features and an \textbf{exercise}.
  \end{block}
\end{frame}
